\chapter{Requirement Analysis}
\label{chap:requirement-analysis}

\section{Stakeholder Analysis}
\label{section:stakeholder-analysis}

Stakeholders are individuals, groups, or entities that have an interest in, or
are affected by, the outcomes of a project. The plugin involves
five distinct stakeholder groups, each with different relationships to the
system and different expectations of its behaviour.

\textbf{Game Developer (Direct User).} The game developer is the primary
technical user of the plugin. This stakeholder integrates the plugin into an
Unreal Engine 5 project, defines which blend shape and morph target parameters
the plugin may control, and selects the LLM operating mode — local inference via
llama.cpp or cloud fallback through a supported provider. The game developer is
responsible for exposing the plugin's character creation widget within the game's
user interface and for configuring any API credentials required for cloud
operation. Their principal concern is that the plugin integrates cleanly as a
drop-in component without requiring modifications to existing skeletal mesh
assets or animation blueprints.

\textbf{Player (Indirect User).} The player interacts with the system
exclusively through the natural-language character creation interface presented
at runtime. This stakeholder has no awareness of the underlying LLM or parameter
mapping mechanism; they expect to describe a character in plain text and receive
a visually faithful result on the skeletal mesh. The player's primary concern is
expressiveness — that the system correctly interprets descriptive terms relating
to facial structure, proportions, and appearance — and responsiveness, in that
the generation latency does not disrupt the flow of gameplay.

\textbf{Project Team.} The project is developed by two core developers
responsible for the plugin architecture, LLM integration, and parameter mapping
pipeline. A third team member contributes to dataset curation and evaluation,
ensuring that the prompt--parameter mapping is validated against a representative
range of character descriptions and skeletal mesh configurations. The team's
collective interest is in delivering a system that is technically sound,
reproducible, and suitable for academic assessment.

\textbf{Academic Advisors.} The primary academic advisor provides guidance on
research methodology, system design, and evaluation criteria. A co-advisor is
yet to be confirmed. The advisors' interest lies in ensuring the project meets
the academic rigour required for a senior project submission, including clear
problem framing, traceable requirements, and a principled evaluation of the
system's performance against its stated objectives.


\section{User Stories and Use Cases}
\label{section:user-stories-use-cases}
<TIP: For each stakeholder, write user stories that capture what
they need from the system: "As a [stakeholder], I want to [action] so that [benefit]."
For your most important user stories, expand them into use cases.
A use case is a detailed description of how a system
interacts with an external entity (such as a user or another system) to
accomplish a specific goal. Use cases provide a high-level view of the
functionality of a system and help in capturing and documenting its
requirements from the perspective of end users.

Make sure to use the
appropriate type of use case for each scenario (brief, casual, and fully-dressed
use case). />

\section{Use Case Diagrams}
\label{section:use-case-diagrams}
<TIP: Include a UML Use Case Diagram that shows all actors and
their relationships to the use cases you described.
Also provide the written description for each use case. />

\section{User Interface Design}
\label{section:user-interface-design}
<TIP: Put the initial design of your application here. You can
showcase a detailed design of a specific page or a sitemap of your application.
See an example below./>

\begin{figure}[h]
    \centering
    \includegraphics[width=0.8\textwidth]{examples/user-interface-design.png}
    \caption{User Interface Design}
\end{figure}